\documentclass[11pt]{article}
\usepackage[margin=1in]{geometry}
\usepackage{amsmath,amssymb,amsthm}
\usepackage{hyperref}
\newtheorem{theorem}{Theorem}
\newtheorem{lemma}[theorem]{Lemma}
\newtheorem{proposition}[theorem]{Proposition}
\newtheorem{corollary}[theorem]{Corollary}
\theoremstyle{definition}
\newtheorem{definition}[theorem]{Definition}
\theoremstyle{remark}
\newtheorem{remark}[theorem]{Remark}
\newcommand{\F}{\mathbb{F}}
\newcommand{\Z}{\mathrm{Z}}
\newcommand{\slt}{\mathfrak{sl}_2}
\newcommand{\sln}{\mathfrak{sl}_n}
\newcommand{\g}{\mathfrak{g}}
\newcommand{\kdim}{\operatorname{Kdim}}
\newcommand{\End}{\operatorname{End}}

\title{Disproof of Conjecture 00000001141:\\
the centre of $u(\mathfrak{sl}_n)$ has Krull dimension $0$, not $n-1$}
\author{jilint777}
\date{2026-10-04}

\begin{document}
\sloppy
\maketitle

\begin{abstract}
Conjecture 00000001141 asserts that the Krull dimension of the centre of the restricted
enveloping algebra $u(\sln)$ is $n-1$, and that for every reductive $\g$ the centre of $u(\g)$ is
spanned by the Weyl-invariant polynomial ring of the $p$-centre, with dimension equal to the
rank. The first clause is false for every $n\ge 2$ and every prime $p$. The algebra $u(\g)$ has
finite dimension $p^{\dim\g}$, so its centre is a finite-dimensional commutative algebra. Every
prime ideal of such an algebra is maximal, so its Krull dimension is $0\neq n-1$. We make this
concrete for $n=2$, $p=3$. The algebra $u(\slt)$ has dimension $27$. Its centre $\Z$ has
dimension $4$ and $81$ elements, with $\Z\cong\F_3\times\F_3[x,y]/(x,y)^2$. It has exactly two
prime ideals, both maximal. The Lean~4 development (core library only) builds a model of
$u(\slt)$ over $\F_3$ from scratch. It proves the defining relations and that the model has
exactly $3^{27}$ elements, and it proves from first principles that its centre has Krull
dimension $0$; in particular this dimension is not $2-1$. It also shows that the centre is not one-dimensional, so the
``dimension equals the rank'' part of the second clause fails for $\g=\slt$ as well. A
standard-library Python script recomputes everything for $p=2,3,5$ by a different method. We
also discuss the other readings. For the full enveloping algebra $U(\sln)$ in characteristic
$p$, the Krull dimension of the centre is $n^2-1$, which again is not $n-1$. Only the
characteristic-$0$ statement for $U(\sln)$ gives $n-1$ (Harish-Chandra), and that statement is
true but is not the one made.
\end{abstract}

\section{The statement}

\begin{quote}
\emph{Definition: The restricted enveloping algebra $u(\g)$ of a restricted Lie algebra $\g$.
Conjecture: The Krull dimension of the center of $u(\sln)$ is $n-1$ (generated by the $p$-centers
of the various $\slt$-sections), and for every reductive $\g$ the center of $u(\g)$ is spanned by
the Weyl-invariant polynomial ring of the $p$-center, with dimension always equal to the rank.}
\end{quote}
The Chinese version says the same thing. The conjecture is a conjunction, so it is enough to refute the clause
\[
\textbf{(C1)}\qquad \kdim \Z\bigl(u(\sln)\bigr)=n-1 \qquad(\text{for every } n),
\]
and we refute it at $n=2$: the claim is $\kdim\Z(u(\slt))=1$. Section~\ref{sec:second} shows that
the second clause fails for $\g=\slt$ as well.

\section{Background}

Let $k$ be a field of characteristic $p>0$. A \emph{restricted Lie algebra} is a Lie algebra $\g$
over $k$ with a $p$-map $x\mapsto x^{[p]}$ such that $\operatorname{ad}(x^{[p]})=(\operatorname{ad}x)^p$ (plus
the usual axioms). The \emph{restricted enveloping algebra} is
\[
u(\g)=U(\g)\big/\bigl(x^p-x^{[p]}\;:\;x\in\g\bigr).
\]
It suffices to take $x$ in a basis of $\g$. Jacobson's PBW theorem for $u(\g)$ states that if
$x_1,\dots,x_d$ is a basis of $\g$, the ordered monomials $x_1^{a_1}\cdots x_d^{a_d}$ with
$0\le a_i<p$ form a basis of $u(\g)$. In particular $\dim_k u(\g)=p^{\dim\g}$ (see Jacobson,
\emph{Lie Algebras}, Ch.~V, or Strade--Farnsteiner, \emph{Modular Lie Algebras and their
Representations}, Ch.~2).

For $\sln$ the $p$-map is the matrix $p$-th power. On the basis
$e=\begin{pmatrix}0&1\\0&0\end{pmatrix}$, $h=\begin{pmatrix}1&0\\0&-1\end{pmatrix}$,
$f=\begin{pmatrix}0&0\\1&0\end{pmatrix}$ of $\slt$ it is
$e^{[p]}=0$, $f^{[p]}=0$, $h^{[p]}=h$. For $p$ odd the centre of $\slt$ is $0$, so this $p$-map is
the only one. Hence
\[
u(\slt)=k\langle e,h,f\rangle\big/\bigl([h,e]=2e,\ [h,f]=-2f,\ [e,f]=h,\ e^p=0,\ f^p=0,\ h^p=h\bigr).
\]

The \emph{Krull dimension} of a commutative ring $R$ is the supremum of the lengths $k$ of
chains $P_0\subsetneq P_1\subsetneq\cdots\subsetneq P_k$ of prime ideals.

\section{The first clause is false for all \texorpdfstring{$n\ge 2$}{n >= 2}}

\begin{lemma}\label{lem:artin}
Let $Z$ be a nonzero commutative algebra of finite dimension over a field $k$. Every prime ideal
of $Z$ is maximal, and $\kdim Z=0$.
\end{lemma}
\begin{proof}
Let $P$ be prime. Then $D=Z/P$ is a finite-dimensional $k$-algebra without zero divisors. For
$0\neq x\in D$, multiplication by $x$ is an injective $k$-linear endomorphism of $D$, hence
surjective, so $xy=1$ for some $y$. So $D$ is a field and $P$ is maximal. Two maximal ideals are
never strictly contained in one another, so there is no chain of primes of length $\ge 1$. A
maximal ideal exists because $Z\neq 0$, so $\kdim Z=0$.
\end{proof}

\begin{theorem}\label{thm:main}
For every field $k$ of characteristic $p>0$ and every $n\ge 1$, $\kdim\Z(u(\sln))=0$. In
particular $\kdim\Z(u(\sln))\neq n-1$ for all $n\ge 2$.
\end{theorem}
\begin{proof}
$u(\sln)$ has dimension $p^{n^2-1}$, so its centre is a finite-dimensional commutative algebra.
It contains $k\cdot 1\neq 0$. Apply Lemma~\ref{lem:artin}.
\end{proof}

The same argument applies to every finite-dimensional algebra. This covers $u(\g)$ for any
restricted $\g$ and the reduced enveloping algebras $u_\chi(\g)$. When $k=\F_q$ is finite, the
centre is a finite ring. For finite rings the proof below is the one formalized in Lean. It
uses only that powers of an element repeat.

\begin{lemma}[finite rings]\label{lem:finite}
Let $R$ be a finite commutative ring and $P\subseteq Q$ prime ideals. Then $P=Q$.
\end{lemma}
\begin{proof}
Let $x\in Q$. By the pigeonhole principle $x^i=x^j$ for some $i<j$. Put $d=j-i$; then
$x^{m+d}=x^m$ for all $m\ge i$, so $y=x^n$ with $n=d(i+1)$ satisfies $n\ge 1$ and $y^2=y$. Hence
$y(1-y)=0\in P$, so $y\in P$ or $1-y\in P$. If $1-y\in P\subseteq Q$, then since $y\in Q$ also
$1=y+(1-y)\in Q$, which is impossible. So $x^n=y\in P$, and since $P$ is prime, $x\in P$.
\end{proof}

\section{The example \texorpdfstring{$n=2$, $p=3$}{n=2, p=3} in detail}

\subsection{Normal ordering}
\begin{lemma}\label{lem:formulas}
In any associative algebra generated by $e,h,f$ with $[h,e]=2e$, $[h,f]=-2f$, $[e,f]=h$, for
all $a\ge 0$ and every polynomial $q$:
\[
h\,e^a=e^a(h+2a),\qquad q(h)\,e=e\,q(h+2),\qquad f\,q(h)=q(h+2)\,f,\qquad
[e^a,f]=a\,e^{a-1}(h+a-1).
\]
Consequently, for the monomials $m_{abc}=e^ah^bf^c$,
\begin{align*}
e\cdot m_{abc}&=m_{a+1,b,c},\qquad
h\cdot m_{abc}=e^a(h+2a)h^bf^c,\\
f\cdot m_{abc}&=e^a(h+2)^bf^{c+1}-a\,e^{a-1}(h+a-1)h^bf^c .
\end{align*}
\end{lemma}
\begin{proof}
$he=e(h+2)$ is $[h,e]=2e$, and induction on $a$ gives the first formula. The second follows
from it, and the third is the mirror image $fh=(h+2)f$ of $[h,f]=-2f$. For the last,
$[e^{a+1},f]=e[e^a,f]+[e,f]e^a=a\,e^{a}(h+a-1)+he^a=e^a\bigl(a(h+a-1)+h+2a\bigr)
=(a+1)e^a(h+a)$. The displayed products follow.
\end{proof}

\begin{proposition}[PBW spanning]\label{prop:span}
If $B$ is an associative algebra generated by $e,h,f$ satisfying the relations of $u(\slt)$, then
$B$ is spanned by the $p^3$ monomials $m_{abc}$, $0\le a,b,c<p$. Hence $\dim u(\slt)\le p^3$.
\end{proposition}
\begin{proof}
Let $S$ be the span of the $m_{abc}$. It contains $1$. By Lemma~\ref{lem:formulas}, together with
$e^p=f^p=0$ and $h^p=h$ (which reduce every polynomial in $h$ to degree $<p$), $S$ is stable under
left multiplication by $e$, $h$ and $f$. So $S$ is a left ideal containing $1$, and $S=B$.
\end{proof}
(The reverse inequality is Jacobson's theorem. It is also reproved below for $p=3$, because our
model has exactly $3^{27}$ elements.)

\subsection{A concrete model of \texorpdfstring{$u(\mathfrak{sl}_2)$}{u(sl2)} over \texorpdfstring{$\F_3$}{F3}}
Let $V=\F_3^{27}$ with basis $b_i$, $i=9a+3b+c$, standing for $m_{abc}$. Let $E,H,F\in\End(V)$
be the operators given by the formulas of Lemma~\ref{lem:formulas}, i.e.\ left multiplication by
$e,h,f$ on $u(\slt)$. Let $R_E,R_H,R_F$ be the right multiplications. All six are explicit sparse
$27\times27$ matrices over $\F_3$. Let $A\subseteq\End(V)$ be the unital subalgebra generated by
$E,H,F$.

\begin{proposition}\label{prop:model}
\begin{enumerate}
\item $E,H,F$ satisfy $[H,E]=2E$, $[H,F]=-2F$, $[E,F]=H$, $E^3=0$, $F^3=0$, $H^3=H$.
\item The evaluation map $\mathrm{ev}\colon A\to V$, $\varphi\mapsto\varphi(b_0)$, is bijective.
So $|A|=3^{27}$.
\item Consequently $A\cong u(\slt)$ over $\F_3$.
\end{enumerate}
\end{proposition}
\begin{proof}
(1) All maps involved are linear, so it suffices to compare them on the $27$ basis vectors.
This is a finite computation (done by \texttt{decide} in Lean).

(2) Surjectivity: $E^aH^bF^c(b_0)=b_{9a+3b+c}$ for all $a,b,c<3$ (finite computation), and $A$ is
closed under linear combinations. Injectivity: $R_E,R_H,R_F$ commute with $E,H,F$ (finite
computation), hence with all of $A$. Also $b_{9a+3b+c}=R_F^cR_H^bR_E^a(b_0)$ (finite
computation). So for $\varphi\in A$ we get $\varphi(b_i)=R_F^cR_H^bR_E^a(\varphi(b_0))$, and
$\varphi$ is determined by $\varphi(b_0)$.

(3) By (1) and the universal property of $U(\slt)$ and of its quotient $u(\slt)$, there is a
surjective algebra map $u(\slt)\to A$ sending $e,h,f\mapsto E,H,F$. By
Proposition~\ref{prop:span}, $|u(\slt)|\le 3^{27}=|A|$, so the map is bijective.
\end{proof}

\subsection{The centre}
Gaussian elimination over $\F_3$ (in \texttt{verify.py}) shows that $\Z=\Z(u(\slt))$ has dimension
$4$, with basis
\[
1,\qquad z_2=h+h^2+ef,\qquad z_3=2h+2h^2+2ehf+e^2f^2,\qquad z_4=2e^2hf^2+e^2h^2f^2 .
\]
The Casimir element is $C=ef+fe+\tfrac12h^2=ef+fe+2h^2=2z_2$. The products are
$z_2^2=z_3$, $z_2z_3=-z_3$, $z_3^2=z_3$ and $z_2z_4=z_3z_4=z_4^2=0$. So $\varepsilon=z_3$ is an
idempotent, $\Z\varepsilon=\F_3\varepsilon$, and $\Z(1-\varepsilon)$ is spanned by $1-\varepsilon$,
$x=z_2+z_3$ and $y=z_4$, with $x^2=xy=y^2=0$. Therefore
\[
\Z\;\cong\;\F_3\times\F_3[x,y]/(x,y)^2 .
\]
This ring has exactly two prime ideals, $\F_3\varepsilon\oplus(x,y)$ and $\Z(1-\varepsilon)$. Each
has $27$ elements and both are maximal. The first is the kernel of the trivial character
$e,h,f\mapsto0$. There is no chain $P\subsetneq Q$, so $\kdim\Z=0\neq1$. A brute-force search in
\texttt{verify.py} over all $212$ subspaces of $\Z$ confirms that $\Z$ has $14$ ideals and $2$
primes, with no inclusions between the primes. For comparison, the script also computes
$\dim\Z(u(\slt))=5,4,7$ for $p=2,3,5$, with $2,2,3$ prime ideals respectively.

\section{The Lean formalization}
The directory \texttt{lean4/} contains a Lean~4.19.0 project that uses only the core library.
All objects are defined from scratch.
\begin{itemize}
\item \texttt{F3} (the field $\F_3$), \texttt{V} $=\F_3^{27}$ (functions \texttt{Fin 27} $\to$ \texttt{F3}),
basis vectors \texttt{bv}, finite sums, and additive maps. \texttt{ext\_basis}: two additive maps
that agree on the basis are equal.
\item \texttt{opE, opH, opF, opRE, opRH, opRF}: the six sparse operators (data tables).
\texttt{Gen}: the inductive closure of $\mathrm{id},E,H,F$ under $+$, $-$ and composition, i.e.
the algebra $A$.
\item \texttt{relations}: Proposition~\ref{prop:model}(1). \texttt{ev\_bijective}:
Proposition~\ref{prop:model}(2). Part (3), via Proposition~\ref{prop:span} and the universal
property, is argued in this report and not in Lean.
\item \texttt{CRing}: commutative rings, with \texttt{IsIdeal}, \texttt{IsPrime},
\texttt{HasChain k} (a chain of primes $P_0\subsetneq\cdots\subsetneq P_k$), and
\texttt{KrullDimIs d} (\texttt{HasChain d} holds and every chain has length $\le d$, i.e.\ the
supremum is $d$). \texttt{Finite} means there is an injection into a type in which the image is
covered by a list.
\item \texttt{primes\_incomparable} (Lemma~\ref{lem:finite}), \texttt{finite\_krullDim\_zero},
\texttt{finite\_krullDim\_ne}, and
\texttt{clause\_fails\_for\_finite\_rings}: no finite commutative ring has Krull dimension
$n-1$ for $n\ge2$.
\item \texttt{ZU}: the centre $\{\varphi\in A:\varphi\psi=\psi\varphi\ \forall\psi\in A\}$, with
the commutative ring structure \texttt{ZRing}. \texttt{ZRing\_finite}: it is finite, by
injectivity of \texttt{ev}. \texttt{P0\_prime}: the kernel of the trivial character
$\chi(\varphi)=$ coefficient of $b_0$ in $\varphi(b_0)$ is a prime ideal.
\item Main results: \texttt{krullDim\_center : ZRing.KrullDimIs 0},
\texttt{conjecture\_00000001141\_false :} $\neg$\,\texttt{Clause 2}, where
\texttt{Clause n := ZRing.KrullDimIs (n - 1)}, and \texttt{krullDim\_center\_ne}: no
$d\ge1$ is the Krull dimension.
\item \texttt{center\_not\_rank\_dim}: the Casimir operator $EF+FE+2H^2$ is central (checked on
generators, then by induction), and $\Z(A)$ is not of the form $\F_3\cdot z$, so
$\dim_{\F_3}\Z\ge2>1=\operatorname{rank}\slt$.
\end{itemize}
The build takes under a minute. The \texttt{\#print axioms} report shows only \texttt{propext},
\texttt{Quot.sound} and (for the pigeonhole step, which uses \texttt{by\_cases}) \texttt{Classical.choice}.
There is no \texttt{sorry} and no \texttt{native\_decide}. The operator tables are data. Their meaning is fixed by the
theorems \texttt{relations} and \texttt{ev\_bijective}, and \texttt{verify.py} re-derives them
independently by word rewriting.

\paragraph{Non-vacuity.} \texttt{HasChain 0} holds through the explicit prime \texttt{P0}, so the
Krull dimension is exactly $0$ rather than undefined. The centre is a genuine subring of $A$,
with $\dim_{\F_3}\Z\ge2$ shown in Lean and $=4$ by \texttt{verify.py}. The relations show that
$A$ is a model of $u(\slt)$, not of some smaller quotient.

\section{Readings of the statement}\label{sec:readings}
\begin{enumerate}
\item \textbf{Literal reading} ($u(\sln)$, any field of characteristic $p$, any $p$). Here
$\kdim\Z=0$ by Theorem~\ref{thm:main}, while the claim is $n-1$. The claim fails for every
$n\ge 2$; it holds trivially for $n=1$, where $\mathfrak{sl}_1=0$. The Lean proof treats $n=2$,
$p=3$, $k=\F_3$. Over an extension $K\supseteq\F_3$ the centre is $\Z\otimes_{\F_3}K$, since the
centre is the kernel of a linear map, and it is still finite-dimensional.
\item \textbf{Vector-space dimension instead of Krull dimension.} $\dim\Z(u(\slt))=5,4,7$ for
$p=2,3,5$, never $1$. Lean proves $\dim\ge2$ for $p=3$. So the clause fails under this reading
too.
\item \textbf{$U(\sln)$ in characteristic $p$} (in case ``$u$'' was meant as $U$). The $p$-centre
$\Z_p=k[x^p-x^{[p]}:x\in\text{basis}]$ is a polynomial ring in $n^2-1$ variables, and $U(\sln)$
is a free $\Z_p$-module of rank $p^{n^2-1}$ (Zassenhaus 1954; see Jantzen, \emph{Representations
of Lie algebras in prime characteristic}, NATO ASI Ser.~C 514, 1998). So $\Z_p\subseteq\Z(U)$ and
$\Z(U)$ is a finitely generated $\Z_p$-module (a submodule of a finite module over a Noetherian
ring). Hence $\Z(U)$ is integral over $\Z_p$, and $\kdim\Z(U)=\kdim\Z_p=n^2-1\neq n-1$ for
$n\ge2$. For $n=2$ the value is $3$. \texttt{verify.py} checks that $e^p,f^p,h^p-h$ are central in
$U(\slt)$ for $p=2,3,5$. This reading is in the report only.
\item \textbf{Characteristic $0$, $U(\sln)$.} By Harish-Chandra and Chevalley,
$\Z(U(\sln))\cong S(\mathfrak h)^{W}$ is a polynomial ring in $n-1$ variables, so its Krull
dimension is $n-1$. This true statement is probably what the conjecture garbles. But the
conjecture explicitly concerns the \emph{restricted} enveloping algebra and the $p$-centre,
which exist only in characteristic $p$. The statement must be judged as written, and as written
it is false.
\end{enumerate}

\section{The second clause}\label{sec:second}
For reductive $\g$ the clause asserts that $\Z(u(\g))$ is spanned by ``the Weyl-invariant
polynomial ring of the $p$-centre'', with dimension equal to the rank. In $u(\g)$ the
$p$-centre generators $x^p-x^{[p]}$ are $0$, so the $p$-centre maps onto $k\cdot 1$, and
polynomials in it span only the scalars. For $\g=\slt$ (reductive, rank $1$) the centre contains
the non-scalar Casimir element. Its dimension is $4\neq1$ for $p=3$ ($\ge2$ in Lean), and its
Krull dimension is $0\neq1$. So the second clause is false as well, under either meaning of
``dimension''. Refuting one clause is already enough.

\section{Independent verification}
\texttt{verify.py} (standard library) rebuilds $u(\slt)$ for $p=2,3,5$ by rewriting words in
$e,h,f$ with the defining relations. It then checks:
\begin{itemize}
\item associativity (all triples for $p\le3$; for $p=5$, $(gx)y=g(xy)$ for generators $g$);
\item the relations, and that $m_{abc}=e^ah^bf^c$;
\item the centre: dimension, closure and commutativity, and that the Casimir is central;
\item the prime ideals: brute force for $p=2,3$, and nilradical plus Frobenius for $p=5$;
\item that the $p$-centre is central in $U(\slt)$;
\item that the six Lean tables equal the rewriting result.
\end{itemize}
All checks pass (see \texttt{verification.txt}).

\section{Conclusion}
$u(\sln)$ is finite-dimensional, so its centre is Artinian and has Krull dimension $0$. The
claimed value $n-1$ is wrong for every $n\ge2$. For $n=2$, $p=3$ this is verified in Lean, down
to an explicit model of $u(\slt)$ and a from-scratch notion of Krull dimension. Conjecture
00000001141 is false.

\end{document}
